\documentclass[10pt,a4paper]{amsart}
\usepackage[margin=19mm]{geometry}
\usepackage{amsmath,amssymb,mathtools}
\usepackage[scr=rsfs,cal=euler]{mathalfa}
\usepackage{xfrac}
\usepackage[unicode,hidelinks,bookmarks=false]{hyperref}
\newcommand{\ie}{i.e.\ }
\newcommand{\cf}{cf.\ }
\newcommand{\resp}{respectively\ }
\newcommand{\Subsection}{\subsection}
\providecommand{\nobreakdash}{\nobreak\hbox{-}}
\setlength{\emergencystretch}{2em}
\multlinegap0pt
\usepackage{bm}
\usepackage{bbm}
\usepackage{dsfont}
\usepackage{xspace}
\usepackage{cancel}
\usepackage{mathtools}
\usepackage{amsthm}
\makeatletter
\@ifundefined{theorem}{\newtheorem{theorem}{Theorem}}{}
\@ifundefined{lemma}{\newtheorem{lemma}[theorem]{Lemma}}{}
\makeatother
\def\mathbb#1{#1}%
\def\mathscr#1{#1}%
\title[Ergodic Theorems for Random Walks in Random Environments]{Ergodic Theorems for Random Walks in Random Environments}
\author{Ayan Ghosh}
\date{Source: arXiv:2601.04161v4 (CC BY 4.0)}
\begin{document}
\maketitle

\noindent\emph{Source and licence.} Ergodic Theorems for Random Walks in Random Environments. Version arXiv:2601.04161v4 (CC BY 4.0), distributed under the Creative Commons Attribution 4.0 International licence, as recorded in the arXiv licence metadata for this exact version and shown on the abstract page https://arxiv.org/abs/2601.04161v4. Full rights notice: licenses/arxiv-2601.04161-CC-BY-4.0.txt.\par\smallskip

\noindent\textit{Editorial scope.} Counted unit: the Lemma whose source label is \texttt{equivalence} in the article (source0.tex lines 827--830, source label \texttt{equivalence}), delivered together with the article's own complete original proof of it (source0.tex lines 831--901, one \texttt{proof} environment of 71 lines). No mathematics is written, repaired or re-derived: statement and proof are the article's own text line for line. Required context retained uncounted: none. Every definition, display and label of those lines is retained unchanged. Omitted from this package and not redistributed: 1--826, 902--1004. Nothing inside a retained excerpt is omitted or altered: every retained body is delivered exactly as the article's lines are written, so no omission receipt of any kind accompanies this package. Citations: the retained excerpts carry no citation command at all, so no bibliography entry of the article is transcribed and no reference is redistributed. Reference adaptation: none was needed and none was made; every reference inside the delivered unit resolves inside it, so no unresolved reference or citation is produced. Printed slips kept literal: no printed word, symbol, hypothesis or number of the counted statement or of its proof was altered. Layout and class: the article's own class is \texttt{aomart} with packages including lmodern, microtype, amsmath, amssymb, amsthm, amsfonts, mathtools, mathrsfs, graphicx, tikz, tikz-cd, babel, csquotes, enumitem; the wrapper is the lane's pinned \texttt{amsart} editorial wrapper at 10pt on A4, which restates the theorem-like environments the retained text needs and adds the article's own macro definitions that the retained text needs (\texttt{\textbackslash mathbb}, \texttt{\textbackslash mathscr}), with extra wrapper packages (bm, bbm, dsfont, xspace, cancel, mathtools). The delivered numbering is the unit's own. Assets: no retained span contains a figure environment, a graphics-inclusion command, a PDF-page inclusion, a table or a TikZ picture; no image file is copied into this package, the package declares no asset dependency and no image file is redistributed with it. Licence: version arXiv:2601.04161v4 of this article is distributed under the Creative Commons Attribution 4.0 International licence, as recorded in the arXiv licence metadata for this exact version and shown on the abstract page https://arxiv.org/abs/2601.04161v4. A full rights notice is delivered at licenses/arxiv-2601.04161-CC-BY-4.0.txt. Characters: every retained excerpt is pure ASCII, so the delivered engine, XeLaTeX with its default Latin Modern text font, reports no missing character.
\begin{lemma}~\label{equivalence}
    Let $\mathbb{Q}$ be a measure on $(\mathscr{S}(d), \mathcal{B}(d)')$ mutually absolutely continuous to
   $\mu$, such that the local process $\{\pi_0(\tau_{X_{\Lambda(\omega)}(m)}\omega)\}_{m \geq 0}$ is stationary. Further assume $\{\pi_0(\tau_{X_{\Lambda(\omega)}(m)}\omega)\}_{m \geq 0}=\{\pi_0(\tau_{X_{\omega}(m)}\omega)\}_{m \geq 0}$ in distribution under $\mathbb{Q}$. Then $\{\pi_0(\tau_{X_{\omega}(m)}\omega)\}_{m \geq 0}$ is stationary and ergodic under $\mathbb{Q}$.
\end{lemma}

\begin{proof}
Stationarity follows trivially from the assumptions. We just have to show the ergodicity of the local process. 

Let $\mathscr{S}'= S(d)^{\mathbb{N}}$ denote all $S(d)$ valued trajectories of the local process $\{\pi_0(\tau_{X_{\omega}(m)}\omega)\}_{m \geq 0}$.

Let $\theta$ denote the canonical shifts or the $n^{th}$ iterate of the process. Denote by $\tilde{\mathcal{B}}$ the product $\sigma$-algebra. Let $A$ be a set which is invariant under $\theta$, i.e., $\theta^{-1} A = A$. 
 
Let $h(\omega) =\tilde{P}_{\omega}(A) $ where $\tilde{P}_\omega$ is the law if the chain starts from \\$\{\omega(x)\}_{x \in \mathbb{Z}^d}$
We must show 
\[
\tilde{P}_{\mathbb{Q}} (A) = \int_{\mathscr{S}(d)} \tilde{P}_{\omega}(A) d\mathbb{Q}(\omega) \in \{0,1\}
\]
Let $P^*_{\omega}$ denote the law of the local process $\{\pi_0(\tau_{X_{\Lambda(\omega)}(m)}\omega)\}_{m \geq 0}$ starting at $\{\omega(x)\}_{x \in \mathbb{Z}^d}$. 

Therefore, it is enough to show, \[
{P}^*_{\mathbb{Q}} (A) = \int_{\mathscr{S}(d)}P^*_{\omega}(A) d\mathbb{Q}(\omega) \in  \{0, 1\}.
\]

Denote $\bar{\omega}(n) = \tau_{X_{\Lambda(\omega)}(n)} \omega$. 

Start by noting that under the canonical filtration sequence $\mathscr{G}_{k}= \sigma(\omega, \bar{\omega}(1),\dots,\bar{\omega}(k-1))$,  
$h(\bar{\omega}(n))$ is a $\tilde{P}_{\mathbb{Q}}$ martingale since 
\[
\mathbb{E}_{\mathbb{Q}} [ \boldsymbol{1}_{A}| \omega, \bar{\omega}(1),\dots, \bar{\omega}({n}) ] = \mathbb{E}_{\mathbb{Q}} [ \boldsymbol{1}_{A}\circ \theta^n | \omega, \bar{\omega}(1),\dots, \bar{\omega}({n}) ] = h(\bar{\omega}(n)) \quad \mathbb{Q}\text{-a.e.}
\]
By the Levy Upwards Theorem,
\begin{equation}\label{eq:marty}
h(\bar{\omega}(n)) \to \boldsymbol{1}_A \qquad (n \to \infty) \quad\mathbb{Q}\text{-a.e.}
\end{equation}

Next we show $h = \boldsymbol{1}_B$($\mathbb{Q}$ \text{a.e.}) for some $B \in \mathcal{B}(d)'$. If not then \\$\mathbb{Q}( h \in [a,b])>0$ for some $0<a<b<1$. 

But by Birkhoff's Ergodic Theorem, 
\[
\frac{1}{n} \sum_{k=0}^{n-1} \boldsymbol{1}\{h(\bar{\omega}(k))\in [a,b]\} \to \Phi= \mathbb{E}_{\mathbb{Q}}[ \boldsymbol{1}\{h(\omega)\in [a,b]\} | \mathcal{I} ]\quad\mathbb{Q}\text{-a.e.}
\]
where $\mathcal{I}$ is the $\sigma$-algebra of invariant events. 

Observe, 
\[
\mathbb{E}_{\mathbb{Q}} [\Phi] = \mathbb{Q}[h \in [a,b]]>0.
\]
But this contradicts \ref{eq:marty} since $\Phi =0 $ ($\mathbb{Q}$ \text{a.e.}).
Now we note that from the martingale property and the Markov property, and the fact that $\{h(\bar{\omega}(n))\}_{n\geq0}$ is a martingale,
\begin{equation}\label{fund}
\boldsymbol{1}_{B}(\omega) 
  = \mathbb{E}_{\mathbb{Q}}\Big[\,\boldsymbol{1}_{B}(\bar{\omega}(1)) \,\big|\, \omega \,\Big]
  = K_d^{\Lambda} \,\boldsymbol{1}_{B}(\omega)
  \quad \mathbb{Q}\text{-a.e.}
\end{equation}
Combining \ref{fund} with ellipticity,
\begin{equation}{\label{ineq}}
\boldsymbol{1}_{B} \geq \boldsymbol{1}_{B} \circ \tau_e \quad (\mu\text{-\text{a.e.} or } \mathbb{Q}\text{-a.e.}) \quad \forall\, e \in {V}_d.
\end{equation}
Iterating \ref{ineq}, we obtain  
\[
\boldsymbol{1}_{B} \geq \boldsymbol{1}_{B} \circ \tau_x \quad (\mu\text{\text{-a.e.} or } \mathbb{Q}\text{\text{-a.e.}}) \quad \forall\, x \in \mathbb{Z}^d.
\]
Since $\mu[ \omega \in B] = \mu [\tau_{x}\omega\in B]$ for every $x \in \mathbb{Z}^d$, 

\begin{equation}{\label{erg}}
 \boldsymbol{1}_{B} = \boldsymbol{1}_{B} \circ \tau_x \quad \mu\text{\text{-a.e.} or}~\mathbb{Q}\text{\text{-a.e.}}
\end{equation}
Consider the set \[
\tilde{B} = \cap_{x \in \mathbb{Z}^d} \tau_{-x}B .
\] 
Observe $(\tau_x)^{-1} \tilde{B} = \tilde{B}$ which shows $\tilde{B}$ is invariant under shifts. 

From \ref{erg} we get that $ \mu(\tilde{B})= \mu(B) = 0$ or $1$ (by ergodicity). From the mutual absolute continuity of $\mathbb{Q}$ and $ \mu$, 
$\mathbb{Q}(B) =P^*_{\mathbb{Q}}(A)=0$ or $1$. This proves ergodicity of the local process in the environment $\mathbb{Q}$.
\end{proof}

\end{document}
